\section{Conclusion}
\label{sec:conclusion}

This work addresses the optimization of contractual coverage for medical
equipment, with the operational goals of reducing the fraction of
uncovered equipment and improving contractual decision-making. The
problem is formalized as a finite-horizon Markov decision process on a
single piece of equipment and solved with the PPO algorithm, trained
entirely in a reproducible simulator built in the absence of direct
access to Philips France's operational data. The contribution is
applicative: we do not propose a new algorithm, but we provide a clean
MDP formulation, a reproducible training pipeline, and a statistically
grounded comparison with two baselines.
On $100$ evaluation episodes drawn from a shared seed pool, the PPO
agent reaches a mean reward of $55.08 \pm 3.62$ ($95\%$ confidence
interval $[54.4, 55.8]$), significantly above the random policy
($-184.91 \pm 12.80$) and significantly below the threshold-based
heuristic ($60.56 \pm 3.10$; Welch's $t$-test $p \ll 0.01$). A
sensitivity sweep over the heuristic's renewal threshold shows that the
reported heuristic sits almost exactly at the peak of its one-parameter
family, and that PPO already matches or outperforms two of the five
tested thresholds.
Taken together, these observations attest that a standard PPO pipeline
recovers a competent coverage policy from scratch, without receiving any
business rule, and that the residual advantage of the heuristic
reflects the structural simplicity of the considered environment: a
one-dimensional rule on the time to expiration is near-optimal here,
and the added value of reinforcement learning is expected to grow with
the complexity and variability of real environments, where no simple
rule can jointly capture dependencies between cost, risk and temporal
dynamics.
These results rely on a deliberately simplified simulator, a
four-dimensional state representation and a training budget of $3{,}000$
episodes, all of which bound the scope of the conclusions.
Natural extensions include the exploitation of real operational data to
enrich the state representation, the extension of the framework to
heterogeneous portfolios subject to global budget constraints, a
sensitivity study of PPO itself to changes in reward weights and
environmental parameters, and ultimately a supervised deployment in
production.
The presented approach demonstrates the viability of driving contractual
decisions through reinforcement learning and opens the way to a
gradual, empirically justified adoption of learned policies
complementary to existing business rules.
